\begin{beginnerstatement}[Kurz erklärt: control.py]

\path{tools/control.py} ist ein Python-Skript und dient als Kommandozentrale
des Templates. Ein Terminal ist ein textbasiertes Fenster, in dem Programme
durch Eingaben gesteuert werden. Dort stellt das Skript eine
\emph{Command Line Interface} (CLI), also eine Kommandozeilenoberfläche,
bereit. Ein Kommando wie \texttt{run} wählt eine Aufgabe; ein Argument liefert
dazu eine zusätzliche Angabe. \emph{Tooling} bezeichnet die Hilfswerkzeuge für
Abläufe wie Installation, Tests und Builds. Ein Exitcode ist eine Zahl, mit
der das Skript Erfolg, einen fachlichen Fehler oder eine falsche Bedienung
meldet. Ein Traceback zeigt bei einem unerwarteten Python-Fehler die Kette der
aufgerufenen Programmstellen; erwartete Bedien- und Profilfehler werden hier
ohne einen solchen technischen Detailbericht ausgegeben.

Der neue Befehl \texttt{quality} koordiniert solche Prüfwerkzeuge. Ein Linter
sucht bestimmte Fehler und Regelverletzungen, ein Formatter prüft die
Schreibweise des Codes. Eine Rule ID ist eine kurze eindeutige Kennung wie
\texttt{CQ001}; damit lässt sich ein Befund trotz wechselnder Erklärung
wiedererkennen. Der Dispatcher nimmt den Befehl an, während die eigentliche
Quality-Logik in eigenen Modulen liegt.

\end{beginnerstatement}
